\section{ER-diagram oversatt til relasjonsskjema}

Relasjonsskjemaet under er utledet direkte fra ER-diagrammet ved
standardreglene for ER-til-relasjonsoversetting.
Primærnøkler er \underline{understreket}, fremmednøkler er \textit{kursive},
og attributter som er begge deler er \underline{\textit{understreket kursive}}.

% Makroer for relasjonsskjema
\newcommand{\pk}[1]{\underline{#1}}
\newcommand{\fk}[1]{\textit{#1}}
\newcommand{\pkfk}[1]{\underline{\textit{#1}}}

\newcommand{\rel}[2]{\noindent\textbf{#1}(#2)}
\newenvironment{fklist}{%
  \begin{itemize}[noitemsep, topsep=1pt, leftmargin=1.5em, label={$\bullet$}]%
}{%
  \end{itemize}\vspace{6pt}%
}

\vspace{4pt}

% ── Sterke entiteter ──────────────────────────────────────────────────────────
\rel{Fasilitet}{\pk{ID}, type, beskrivelse}

\vspace{6pt}

\rel{Senter}{\pk{ID}, navn, gate, nummer, fra, til}

\vspace{6pt}

\rel{Idrettslag}{\pk{ID}, navn, beskrivelse}

\vspace{6pt}

\rel{Profil}{\pk{ID}, type, fornavn, etternavn, epost, telefon}

\vspace{6pt}

\rel{Aktivitet}{\pk{navn}, kategori, beskrivelse}

\vspace{10pt}

% ── Avhengige entiteter ───────────────────────────────────────────────────────
\rel{har\_fasilitet}{\pkfk{senter\_ID}, \pkfk{fasilitet\_ID}}
\begin{fklist}
  \item \fk{senter\_ID} $\rightarrow$ Senter(\pk{ID})
  \item \fk{fasilitet\_ID} $\rightarrow$ Fasilitet(\pk{ID})
\end{fklist}

\rel{Bemanning}{\pkfk{senter\_ID}, \pk{dag}, start, slutt}
\begin{fklist}
  \item \fk{senter\_ID} $\rightarrow$ Senter(\pk{ID})
\end{fklist}

\rel{Sal}{\pkfk{senter\_ID}, \pk{ID}, type, kapasitet}
\begin{fklist}
  \item \fk{senter\_ID} $\rightarrow$ Senter(\pk{ID})
\end{fklist}

\rel{Tredem\o lle}{\pkfk{senter\_ID}, \pkfk{sal\_ID}, \pk{nummer}, produsent, maks\_hastighet, maks\_stigning}
\begin{fklist}
  \item (\fk{senter\_ID}, \fk{sal\_ID}) $\rightarrow$ Sal(\pk{senter\_ID}, \pk{ID})
\end{fklist}

\rel{Sykkel}{\pkfk{senter\_ID}, \pkfk{sal\_ID}, \pk{nummer}, bluetooth}
\begin{fklist}
  \item (\fk{senter\_ID}, \fk{sal\_ID}) $\rightarrow$ Sal(\pk{senter\_ID}, \pk{ID})
\end{fklist}

\rel{Gruppe}{\pkfk{idrettslags\_ID}, \pk{ID}, type, beskrivelse}
\begin{fklist}
  \item \fk{idrettslags\_ID} $\rightarrow$ Idrettslag(\pk{ID})
\end{fklist}

\rel{Prikk}{\pkfk{profil\_ID}, \pk{ID}, dato}
\begin{fklist}
  \item \fk{profil\_ID} $\rightarrow$ Profil(\pk{ID})
\end{fklist}

\rel{er\_medlem}{\pkfk{idrettslags\_ID}, \pkfk{profil\_ID}}
\begin{fklist}
  \item \fk{idrettslags\_ID} $\rightarrow$ Idrettslag(\pk{ID})
  \item \fk{profil\_ID} $\rightarrow$ Profil(\pk{ID})
\end{fklist}

\rel{Idrettslagstime}{\pkfk{senter\_ID}, \pkfk{sal\_ID}, \pk{ID}, type, start, slutt, dato, beskrivelse, \fk{idrettslags\_ID}, \fk{gruppe\_ID}}
\begin{fklist}
  \item (\fk{senter\_ID}, \fk{sal\_ID}) $\rightarrow$ Sal(\pk{senter\_ID}, \pk{ID})
  \item \fk{idrettslags\_ID} $\rightarrow$ Idrettslag(\pk{ID})
  \item (\fk{idrettslags\_ID}, \fk{gruppe\_ID}) $\rightarrow$ Gruppe(\pk{idrettslags\_ID}, \pk{ID})
\end{fklist}

\rel{m\o ter\_til\_idrett}{\pkfk{senter\_ID}, \pkfk{sal\_ID}, \pkfk{idrettslagstime\_ID}, \pkfk{profil\_ID}}
\begin{fklist}
  \item (\fk{senter\_ID}, \fk{sal\_ID}, \fk{idrettslagstime\_ID}) $\rightarrow$ Idrettslagstime(\pk{senter\_ID}, \pk{sal\_ID}, \pk{ID})
  \item \fk{profil\_ID} $\rightarrow$ Profil(\pk{ID})
\end{fklist}

\rel{Gruppeaktivitet}{\pkfk{senter\_ID}, \pkfk{sal\_ID}, \pk{ID}, start, slutt, dato, \fk{aktivitet\_navn}, \fk{instrukt\_ID}}
\begin{fklist}
  \item (\fk{senter\_ID}, \fk{sal\_ID}) $\rightarrow$ Sal(\pk{senter\_ID}, \pk{ID})
  \item \fk{aktivitet\_navn} $\rightarrow$ Aktivitet(\pk{navn})
  \item \fk{instrukt\_ID} $\rightarrow$ Profil(\pk{ID})
\end{fklist}

\rel{m\o ter\_til\_gruppe}{\pkfk{senter\_ID}, \pkfk{sal\_ID}, \pkfk{gruppeaktivitet\_ID}, \pkfk{profil\_ID}, tidspunkt}
\begin{fklist}
  \item (\fk{senter\_ID}, \fk{sal\_ID}, \fk{gruppeaktivitet\_ID}) $\rightarrow$ Gruppeaktivitet(\pk{senter\_ID}, \pk{sal\_ID}, \pk{ID})
  \item \fk{profil\_ID} $\rightarrow$ Profil(\pk{ID})
\end{fklist}

\rel{p{\aa}meldt\_til}{\pkfk{senter\_ID}, \pkfk{sal\_ID}, \pkfk{gruppeaktivitet\_ID}, \pkfk{profil\_ID}, p{\aa}melding\_nummer}
\begin{fklist}
  \item (\fk{senter\_ID}, \fk{sal\_ID}, \fk{gruppeaktivitet\_ID}) $\rightarrow$ Gruppeaktivitet(\pk{senter\_ID}, \pk{sal\_ID}, \pk{ID})
  \item \fk{profil\_ID} $\rightarrow$ Profil(\pk{ID})
\end{fklist}

\vspace{4pt}

\subsection{Normalformanalyse}

Alle tabellene er på \textbf{BCNF} (Boyce–Codd normalform).
En relasjon er på BCNF dersom venstre side i enhver ikke-triviell funksjonell avhengighet er en supernøkkel.
Nedenfor begrunnes dette gruppevis.

\paragraph{Sterke entiteter – Fasilitet, Senter, Idrettslag, Aktivitet.}
Disse har én enkeltattributt som primærnøkkel (\textit{ID}, resp.\ \textit{navn} for Aktivitet).
Alle øvrige attributter bestemmes entydig av primærnøkkelen, og det finnes ingen andre funksjonelle
avhengigheter. Dermed er venstre side i enhver ikke-triviell FD en supernøkkel, og tabellene er på BCNF.
For \textit{Senter} kunne man vurdert \textit{(gate, nummer)} som kandidatnøkkel, men vi antar at
dette ikke holder: et senter kan i prinsippet flytte eller to sentre kan midlertidig befinne seg i
samme bygg (samme gateadresse og nummer). Adressen er dermed ikke garantert unik og regnes ikke som
kandidatnøkkel.

\paragraph{Profil.}
Primærnøkkelen er \textit{ID}. I tillegg er \textit{epost} unik for hver profil
(brukstilfellene identifiserer brukere ved epostadresse), slik at
\textit{epost} $\rightarrow$ \textit{ID, type, fornavn, etternavn, telefon}
også er en ikke-triviell FD. Venstresiden er en kandidatnøkkel og dermed en supernøkkel,
så BCNF er oppfylt.

\paragraph{Svake entiteter – Bemanning, Sal, Tredemølle, Sykkel, Idrettslagstime, Gruppeaktivitet, Gruppe, Prikk.}
ER-til-relasjonsoversettingen av svake entiteter gir en sammensatt primærnøkkel bestående av
eierens primærnøkkel og diskriminatoren. De eneste ikke-trivielle funksjonelle avhengighetene er
av formen \(\text{primærnøkkel} \rightarrow \text{øvrige attributter}\), der venstresiden er
primærnøkkelen og dermed en supernøkkel. BCNF er derfor oppfylt for alle disse tabellene.

\paragraph{M:N-relasjoner uten ekstra attributter – har\_fasilitet, er\_medlem, møter\_til\_idrett.}
Alle attributter inngår i primærnøkkelen. Det finnes ingen ikke-trivielle funksjonelle avhengigheter
utover de trivielle, og tabellene er derfor trivielt på BCNF.

\paragraph{møter\_til\_gruppe.}
Primærnøkkelen er sammensatt av alle fire fremmednøkkelattributtene.
\textit{tidspunkt} avhenger funksjonelt av hele primærnøkkelen, og det finnes
ingen andre ikke-trivielle FD-er. Tabellen er på BCNF.

\paragraph{påmeldt\_til.}
Primærnøkkelen er (\textit{senter\_ID, sal\_ID, gruppeaktivitet\_ID, profil\_ID}).
\textit{påmelding\_nummer} er unikt innenfor én gruppeaktivitet (kølapp-semantikk),
slik at (\textit{senter\_ID, sal\_ID, gruppeaktivitet\_ID, påmelding\_nummer}) $\rightarrow$ \textit{profil\_ID}
også er en ikke-triviell FD. Venstresiden er en kandidatnøkkel og dermed en supernøkkel.
Begge ikke-trivielle FD-er har altså en supernøkkel på venstresiden, og tabellen er på BCNF.
